% ===== BOT 17a: Dao sau muc 17.1 -- 3 khoang trong bang so lieu that =====

\begin{frame}{Vì sao cộng 3 kênh RGB thay vì luminance -- đo trên ảnh thật}
\footnotesize
Trên crop thật \texttt{samples\_drjohnson.png} (panel "gt IMG\_6292", $440{\times}288$px), tính thêm gradient luminance ITU-R BT.709 để đối chứng với công thức Di Zenzo thật:
\[
\boxed{\ \text{grad}_{\text{DiZenzo}}=\sqrt{I_{xx}+I_{yy}},\ \ I_{xx}=\sum_{\text{ch}}I_x^{(\text{ch})2},\ I_{yy}=\sum_{\text{ch}}I_y^{(\text{ch})2}\ }\quad\text{(\texttt{loss\_utils.py:198--208}) vs.}\quad \text{grad}_{\text{lum}}=\sqrt{L_x^2+L_y^2}
\]
\begin{itemize}
\item Quét toàn crop, tìm pixel thật nơi tỉ lệ $\text{grad}_{\text{DiZenzo}}/\text{grad}_{\text{lum}}$ lớn nhất (vùng grad Di Zenzo top-30\%, tránh nhiễu nền phẳng).
\item Rơi đúng viền bình chữa cháy đỏ áp sát cửa gỗ be, tại $(x{=}293,y{=}260)$: $\text{grad}_{\text{lum}}=0.0029$ nhưng $\text{grad}_{\text{DiZenzo}}=0.3227$ -- \textbf{lớn hơn 81.8 lần}.
\item Per-channel tại điểm đó: $|\text{grad}\,R|=0.309$, $|\text{grad}\,G|=0.091$, $|\text{grad}\,B|=0.021$ -- kênh R mang gần hết tín hiệu cạnh, luminance gần như triệt tiêu vì độ sáng hai bên viền xấp xỉ nhau.
\item Bằng chứng số thật (không phải ví dụ lý thuyết) cho câu "luminance bỏ sót cạnh iso-luminant" ở \texttt{loss\_utils.py:172--173}.
\end{itemize}
\begin{center}
\includegraphics[width=0.82\linewidth]{deep17a_01_iso_luminant_dizenzo.png}
\captionof{figure}{\footnotesize Ảnh thật đánh dấu điểm tìm được, zoom cận cảnh, biểu đồ cột 5 đại lượng tại điểm, và ba bản đồ toàn crop (grad luminance / grad Di Zenzo / tỉ lệ).}
\end{center}
\end{frame}


\begin{frame}{Vai trò của $\rho$ ở bước 4: từ ước lượng điểm ảnh giòn tới coherence thật}
\footnotesize
Cùng $I_{xx},I_{xy},I_{yy}$ thô (bước 3, chưa tích hợp) trên crop thật, so sánh ba trạng thái $\rho=0$ (thô), $\rho=1.5$, $\rho=6.0$:
\[
\boxed{\ S_{xx}=G_\rho*I_{xx},\ S_{xy}=G_\rho*I_{xy},\ S_{yy}=G_\rho*I_{yy}\ }\ \text{(\texttt{loss\_utils.py:210--217})},\qquad
\text{coherence}=\Big(\tfrac{\lambda_1-\lambda_2}{\lambda_1+\lambda_2}\Big)^2
\]
\begin{itemize}
\item Coherence trung bình toàn crop \textbf{giảm} theo $\rho$: $\rho{=}0\!:0.884 \to \rho{=}1.5\!:0.726 \to \rho{=}6.0\!:0.590$ -- ngược với trực giác "blur nhiều hơn thì coherent hơn".
\item Lý do: ở $\rho=0$, $I_{xx}$ tại mỗi pixel là ma trận \textbf{hạng-1} theo đúng một hướng gradient điểm ảnh đó $\Rightarrow$ coherence $\approx1$ \emph{giả tạo} ở hầu hết pixel có cạnh -- không phân biệt được cạnh thẳng thật với góc/texture.
\item Chỉ sau khi $G_\rho$ thật sự trung bình hoá không gian, vùng nhiều hướng gần nhau (góc cửa, khung cửa sổ, chân ghế) mới lộ coherence thấp; đoạn cạnh dài đơn hướng (mép cửa, thanh dàn sưởi) vẫn giữ coherence cao, mượt hơn hẳn bản đồ thô đầy nhiễu hạt.
\item Đây là bằng chứng số cho lý do bước 4 \textbf{bắt buộc} phải có $G_\rho$: không có nó, coherence "ảo" ở khắp nơi, không phân biệt cạnh với góc/texture (khớp lý luận đã nêu ở slide "Tensor cấu trúc $J_\rho$").
\end{itemize}
\begin{center}
\includegraphics[width=0.86\linewidth]{deep17a_02_rho_integration_window.png}
\captionof{figure}{\footnotesize $S_{xx},S_{xy},S_{yy}$ và coherence ở $\rho{=}0$/$1.5$/$6.0$ trên cùng crop thật -- $\rho$ lớn hơn "gộp" cạnh lân cận thành mảng liền, lộ đúng vùng góc/texture thay vì nhiễu điểm ảnh giả-coherent.}
\end{center}
\end{frame}


\begin{frame}{Chuẩn hoá batch ở bước 5: bất biến thật với độ sáng, trừ vùng cháy sáng}
\footnotesize
Cùng crop thật, nhân trực tiếp pixel với hệ số độ sáng $k\in\{0.5,1.0,1.5\}$ (clip $[0,1]$ như ảnh 8-bit thật):
\[
\boxed{\ S_{xx},S_{xy},S_{yy}\ \leftarrow\ \dfrac{S_{xx},S_{xy},S_{yy}}{\max_{\text{batch}}(S_{xx}+S_{yy})+10^{-6}}\ }\quad\text{(\texttt{loss\_utils.py:219--227}})
\]
\begin{itemize}
\item \textbf{Trước} chuẩn hoá: $\max(S_{xx}{+}S_{yy})=9.730\ (\times0.5)$, $38.920\ (\times1.0)$, $44.275\ (\times1.5)$ -- lệch rất lớn theo độ sáng.
\item Tỉ lệ $\times0.5$: đo được $9.730/38.920=0.250$, khớp gần tuyệt đối lý thuyết bậc hai $0.5^2=0.250$ (không bão hoà).
\item \textbf{Sau} chuẩn hoá: bản đồ $S_{xx}$ của $\times0.5$ và $\times1.0$ \textbf{giống hệt nhau} -- $\max|\text{diff}|=0.0000$ toàn crop -- vì $k^2$ triệt tiêu đúng đại số ở tử/mẫu tại từng pixel.
\item Với $\times1.5$: đa số pixel vẫn giống hệt, nhưng $\max|\text{diff}|=0.3032$ tại vùng cửa sổ/khung cửa bị \textbf{bão hoà} $[0,1]$ -- tỉ lệ max đo được chỉ $1.138$ thay vì $1.5^2{=}2.25$ lý thuyết, đúng dấu hiệu giả định tuyến tính bị phá vỡ khi ảnh over-expose.
\item Kết luận thực nghiệm: ngưỡng $\eta$ (mục 17.3) bất biến \emph{chính xác} với độ sáng scene ở vùng không bão hoà -- một giới hạn thật đã đo được, không phải giả định suông.
\end{itemize}
\begin{center}
\includegraphics[width=0.62\linewidth]{deep17a_03_brightness_normalization.png}
\captionof{figure}{\footnotesize Trên: crop thật ở 3 mức sáng. Giữa: $S_{xx}{+}S_{yy}$ trước chuẩn hoá (lệch $\sim$4.5 lần). Dưới: $S_{xx}$ sau chuẩn hoá (gần như giống hệt) và hai bản đồ hiệu số.}
\end{center}
\end{frame}
